\chapter{Creatinine}

\section*{What Is This Test?}
The creatinine blood test measures the concentration of creatinine in the serum or plasma. Creatinine is a waste product formed from the non-enzymatic, spontaneous dehydration and cyclization of creatine and creatine phosphate in skeletal muscle. Creatine is synthesized in the liver from arginine, glycine, and methionine, then transported to muscle where it is phosphorylated to creatine phosphate (phosphocreatine), which serves as a high-energy phosphate reservoir for rapid ATP regeneration during muscle contraction. Approximately 1--2\% of total muscle creatine is converted to creatinine daily and released into the blood at a relatively constant rate proportional to total muscle mass. Creatinine is freely filtered at the glomerulus, is not reabsorbed, and undergoes only minor active tubular secretion (approximately 10--15\% of total excretion, which increases as GFR declines). These properties make serum creatinine the most widely used endogenous biomarker for estimating glomerular filtration rate (GFR). Serum creatinine concentration reflects the balance between muscle production (which is constant in a given individual) and renal excretion. A doubling of serum creatinine corresponds to approximately a 50\% reduction in GFR. Creatinine is a component of the basic metabolic panel (BMP) and comprehensive metabolic panel (CMP). It is used to calculate the estimated GFR (eGFR) using equations such as the CKD-EPI (Chronic Kidney Disease Epidemiology Collaboration) 2021 equation, which incorporates serum creatinine, age, sex, and (in earlier versions) race \cite{inker2021newequations}.

\section*{Alternative Names}
\begin{itemize}
  \item Serum Creatinine
  \item Blood Creatinine
  \item Creat
  \item Cr
  \item Plasma Creatinine
  \item Creatinine Clearance (urine-based test)
  \item eGFR
\end{itemize}
\section*{Principle}
Serum creatinine is most commonly measured by the Jaffe reaction (alkaline picrate method), first described in 1886 \cite{tietz2023textbook}. Creatinine reacts with picric acid in alkaline solution to form a yellow-orange colored complex (Janovsky complex) measured spectrophotometrically at 500--520 nm. The rate of color development is proportional to creatinine concentration. The kinetic Jaffe method measures the rate of color formation over a specific time interval, reducing interference from non-creatinine chromogens (such as glucose, ascorbic acid, acetone, acetoacetate, pyruvate, uric acid, and cephalosporins) that react more slowly with alkaline picrate. However, the Jaffe method still overestimates creatinine by 10--20\% compared to reference methods due to persistent chromogen interference. Enzymatic methods using creatininase (creatinine amidohydrolase), creatinase, and sarcosine oxidase coupled to a peroxidase indicator reaction are more specific and accurate. The sequence: Creatinine + H\textsubscript{2}O $\rightarrow$ Creatine (creatininase); Creatine + H\textsubscript{2}O $\rightarrow$ Sarcosine + Urea (creatinase); Sarcosine + O\textsubscript{2} + H\textsubscript{2}O $\rightarrow$ Glycine + Formaldehyde + H\textsubscript{2}O\textsubscript{2} (sarcosine oxidase); H\textsubscript{2}O\textsubscript{2} + chromogen $\rightarrow$ colored product (peroxidase). The enzymatic method has largely replaced the Jaffe method in modern analyzers and produces results that are traceable to the reference method (isotope-dilution mass spectrometry, IDMS). IDMS-standardized creatinine values are approximately 5--10\% lower than non-standardized values. The eGFR equation (CKD-EPI) is calibrated for IDMS-traceable creatinine methods. This test is performed on serum separator tubes using spectrophotometric/enzymatic methods, not by hematology analyzers.

\section*{History}
Creatinine was first identified by Michel Eugène Chevreul in 1832 and named from the Greek ``kreas'' (flesh). Max Jaffe described the alkaline picrate reaction for creatinine in 1886, which remained the standard clinical method for over 130 years. The kinetic modification to improve specificity was developed in the 1960s. The enzymatic creatinine method was introduced in the 1980s. The relationship between serum creatinine and GFR was established by Homer Smith in the 1950s. The Cockcroft-Gault equation for estimating creatinine clearance was published in 1976. The MDRD (Modification of Diet in Renal Disease) study equation was published in 1999. The CKD-EPI equation, which is more accurate than MDRD especially at higher GFR levels, was published in 2009 and updated in 2021 to remove race as a variable \cite{levey2009newequation}. IDMS standardization of creatinine assays was implemented globally in the 2000s.

\section*{Why This Test Is Ordered}
\begin{itemize}
  \item Estimation of glomerular filtration rate (GFR): eGFR is calculated from serum creatinine, age, and sex using the CKD-EPI equation and is the standard for assessing kidney function
  \item Screening for and staging of chronic kidney disease (CKD): KDIGO defines CKD as eGFR <60 mL/min/1.73 m\textsuperscript{2} or markers of kidney damage persisting >3 months \cite{kdigo2013kdigo}
  \item Diagnosis of acute kidney injury (AKI): KDIGO defines AKI as an increase in serum creatinine of $\geq$0.3 mg/dL within 48 hours, or $\geq$1.5 times baseline within 7 days \cite{bellomo2012acute}
  \item Monitoring of renal function during nephrotoxic drug therapy: aminoglycosides, vancomycin, cisplatin, amphotericin B, NSAIDs, ACE inhibitors, ARBs, calcineurin inhibitors, contrast dye
  \item Drug dosing for renally excreted medications: numerous drugs require dose adjustment based on creatinine clearance or eGFR (e.g., enoxaparin, metformin, gabapentin, many antibiotics)
  \item Preoperative assessment as part of the BMP/CMP
  \item Monitoring of kidney transplant function and early detection of rejection or calcineurin inhibitor toxicity
  \item Evaluation of patients with hypertension, diabetes, cardiovascular disease, or proteinuria for concomitant renal disease
  \item Assessment of volume status: creatinine is less affected by hydration status than BUN, so the BUN:Cr ratio is more informative than either alone
  \item Screening for contrast-induced nephropathy risk: eGFR <30 precludes certain contrast procedures; eGFR 30--60 requires hydration protocols
\end{itemize}

\section*{Primary vs Secondary Test}
Serum creatinine with eGFR is a primary screening test and is arguably the single most important laboratory test for renal function assessment. It is part of every BMP and CMP and is ordered millions of times daily worldwide. It is the cornerstone of CKD screening, staging, and monitoring.

\section*{How This Test Is Performed}
A healthcare professional draws blood from a vein into a serum separator tube (gold-top) or lithium heparin tube (green-top). Approximately 3--5 mL of blood is collected. Standard venipuncture technique is used. The sample is centrifuged and creatinine measured on an automated chemistry analyzer. Results are available within 30--60 minutes \cite{tierney2023current}. The laboratory information system automatically calculates the eGFR using the CKD-EPI equation and reports it alongside the creatinine value. No special handling is required. For point-of-care creatinine testing, fingerstick whole blood is used with handheld analyzers, though results may differ slightly from central laboratory values.

\section*{What Preparation Is Needed}
No fasting is required for serum creatinine. Important considerations for accurate interpretation: (1) Strenuous exercise within 24 hours can transiently elevate creatinine from increased muscle creatine phosphate breakdown (minimal in most people but may be noticeable after extreme exertion or in athletes). (2) Recent consumption of a large meat meal (especially cooked/grilled meat) can transiently increase creatinine by 0.2--0.3 mg/dL within 2--4 hours from absorption of preformed creatinine in cooked muscle. Fasting morning samples avoid this. (3) Certain medications affect creatinine measurement: cephalosporins (especially cefoxitin, cephalothin) interfere with the Jaffe reaction, causing falsely elevated creatinine by up to 1--2 mg/dL; enzymatic methods are not affected. Flucytosine, methyldopa, and levodopa also interfere with Jaffe methods. (4) Medications that truly increase creatinine by reducing GFR without true nephrotoxicity: ACE inhibitors/ARBs (hemodynamic effect---10--20\% rise from reduced glomerular capillary pressure, acceptable up to 30\% rise), trimethoprim (blocks tubular creatinine secretion, raising creatinine by 0.2--0.4 mg/dL without reducing GFR), cimetidine (also blocks tubular secretion, similar effect but now rarely used), dolutegravir (similar to trimethoprim). (5) Bilirubin >15 mg/dL interferes with some enzymatic and Jaffe creatinine methods (negative interference). (6) Ketones (acetoacetate) interfere with the Jaffe reaction in DKA, causing pseudohypercreatininemia. (7) Hemolysis and lipemia can interfere spectrophotometrically \cite{wallach2022interpretation}.

\section*{Contraindications and Precautions}
\begin{itemize}
  \item No absolute contraindications. Critical considerations for creatinine interpretation:
  \item (1) Serum creatinine is proportional to muscle mass: a ``normal'' creatinine of 1.0 mg/dL in a frail, elderly, cachectic woman with muscle wasting may represent an eGFR <30 mL/min. Conversely, a creatinine of 1.5 mg/dL in a bodybuilder may represent an eGFR >90 mL/min. Always assess muscle mass clinically when interpreting creatinine and eGFR.
  \item (2) False negatives for renal failure: elderly malnourished patients, amputees, cirrhotic patients, patients with muscular dystrophy, and those with spinal cord injury may have ``normal'' creatinine despite severely reduced GFR due to low muscle mass. Cystatin C is a useful alternative in these patients.
  \item (3) False positives for renal failure: high muscle mass (bodybuilders, large muscular males), high meat intake, creatine supplements (elevate creatinine from increased creatine turnover), rhabdomyolysis (increased creatinine from muscle breakdown), and drugs blocking tubular secretion (trimethoprim, cimetidine).
  \item (4) Pregnancy: GFR increases by 40--50\% in pregnancy, so serum creatinine normally decreases to 0.4--0.6 mg/dL. A creatinine >0.8 mg/dL in pregnancy is abnormal and warrants investigation.
  \item (5) The eGFR is less accurate at GFR >60 mL/min/1.73 m\textsuperscript{2} and should be reported as ``>60'' rather than a specific number per KDIGO guidelines (though many labs report specific values up to 90).
  \item (6) eGFR equations assume stable renal function and are less accurate in AKI where creatinine is changing rapidly.
\end{itemize}
\section*{Risks}
Minimal risk from standard venipuncture. Mild pain or bruising at puncture site resolves in 1--2 days. Rare: hematoma, infection, nerve injury, vasovagal syncope. Blood volume drawn (3--5 mL) is negligible.

\section*{What Happens After the Test}
Results are available within 30--60 minutes. Serum creatinine is never interpreted alone---the eGFR is calculated and reported. The KDIGO CKD classification is: Stage 1: eGFR $\geq$90 with kidney damage markers (albuminuria, hematuria, structural abnormalities); Stage 2: eGFR 60--89 with kidney damage; Stage 3a: eGFR 45--59; Stage 3b: eGFR 30--44; Stage 4: eGFR 15--29; Stage 5: eGFR <15 or on dialysis. Critical values: acute rise in creatinine $\geq$0.3 mg/dL over 48 hours or rise to $\geq$1.5 times baseline (KDIGO AKI Stage 1). AKI Stage 2: $\geq$2.0--2.9 times baseline; Stage 3: $\geq$3.0 times baseline or creatinine $\geq$4.0 mg/dL with acute rise $\geq$0.3 or initiation of RRT. The rate of creatinine rise in AKI is clinically significant: a slow rise (over days) may indicate resolving ATN; a rapid rise (over hours) suggests severe acute insult. Serial creatinine monitoring every 12--24 hours is standard for AKI management and following response to treatment. For CKD, creatinine and eGFR should be monitored every 3--12 months depending on stage, progression rate, and complications.

\section*{Understanding Results}

\subsection*{Below Normal}
\begin{itemize}
  \item \textbf{Low creatinine (<0.5 mg/dL in adult females, <0.7 mg/dL in adult males):} Not clinically significant in the absence of renal pathology and generally does not indicate supernormal renal function. Most commonly reflects low muscle mass.
  \item \textbf{Reduced muscle mass:} Elderly, frail, cachectic patients with severe malnutrition, anorexia nervosa, prolonged critical illness, amputees, paraplegia/quadriplegia from spinal cord injury, muscular dystrophy (Duchenne, Becker, myotonic dystrophy). In these patients, creatinine may be deceptively normal even with advanced CKD.
  \item \textbf{Pregnancy:} Physiological increase in GFR by 40--50\% lowers creatinine to 0.4--0.6 mg/dL by the second trimester. Creatinine >0.8 mg/dL in pregnancy is abnormal.
  \item \textbf{Advanced cirrhosis:} Reduced creatine synthesis in the liver plus reduced muscle mass.
  \item \textbf{Overaggressive IV fluids:} Dilutional effect.
  \item \textbf{Steroid or neuromuscular blocking agent-induced:} In ICU patients, prolonged neuromuscular blockade and steroid myopathy decrease muscle mass and lower creatinine over time.
\end{itemize}

\subsection*{Above Normal}
\begin{itemize}
  \item \textbf{Acute kidney injury (AKI):} Prerenal (volume depletion, CHF, cirrhosis, renal artery stenosis, NSAIDs, ACE-I/ARBs, hepatorenal syndrome); Intrinsic renal: ATN (prolonged prerenal, ischemia, nephrotoxins---aminoglycosides, cisplatin, amphotericin B, contrast dye, myoglobin, hemoglobin), acute glomerulonephritis (post-streptococcal, anti-GBM, ANCA vasculitis, lupus nephritis), acute interstitial nephritis (drugs: NSAIDs, penicillins, cephalosporins, sulfonamides, PPIs), atheroembolic disease, thrombotic microangiopathy (HUS/TTP), myeloma cast nephropathy; Postrenal: BPH, prostate cancer, bilateral ureteral obstruction, neurogenic bladder.
  \item \textbf{Chronic kidney disease (CKD):} Diabetic nephropathy (leading cause worldwide, 40--50\% of ESRD cases), hypertensive nephrosclerosis (second leading cause), chronic glomerulonephritis (IgA nephropathy, FSGS, membranous nephropathy, membranoproliferative GN), polycystic kidney disease, chronic interstitial nephritis (analgesic nephropathy, lithium, heavy metals, autoimmune), obstructive uropathy, recurrent pyelonephritis/vesicoureteral reflux, HIV-associated nephropathy, sickle cell nephropathy, amyloidosis, light chain deposition disease.
  \item \textbf{Non-renal causes of elevated creatinine (without reduced GFR):} High muscle mass (bodybuilding, large muscular male---may have creatinine of 1.3--1.5 mg/dL with normal GFR); heavy meat meal within 2--4 hours; creatine/creatinine supplements; rhabdomyolysis (increased creatinine from muscle breakdown, but GFR is also usually reduced from ATN); drugs blocking tubular secretion: trimethoprim, cimetidine, dolutegravir, bictegravir, rilpivirine (raise creatinine 0.2--0.4 mg/dL without affecting GFR---check cystatin C if concern).
  \item \textbf{An acute serum creatinine rise of >0.3 mg/dL in 48 hours triggers AKI alert per KDIGO and should never be ignored.}
\end{itemize}

\section*{Normal Results}
\begin{itemize}
  \item \textbf{Adult males:} 0.7--1.3 mg/dL (62--115 $\mu$mol/L) \cite{fischbach2023manual}
  \item \textbf{Adult females:} 0.5--1.1 mg/dL (44--97 $\mu$mol/L)
  \item \textbf{Children (1--16 years):} 0.3--1.0 mg/dL (age and sex-dependent)
  \item \textbf{Infants (1--12 months):} 0.2--0.4 mg/dL
  \item \textbf{Newborns (0--7 days):} 0.3--1.0 mg/dL (reflects maternal creatinine initially, falls over first week)
  \item \textbf{Pregnancy:} 0.4--0.6 mg/dL
  \item \textbf{eGFR (CKD-EPI):} $\geq$90 mL/min/1.73 m\textsuperscript{2} (normal)
  \item \textbf{CKD stage definitions:} Stage 1 (eGFR $\geq$90 with damage), Stage 2 (60--89 with damage), Stage 3a (45--59), Stage 3b (30--44), Stage 4 (15--29), Stage 5 (<15 or dialysis)
  \item \textbf{Conversion:} 1 mg/dL = 88.4 $\mu$mol/L; 1 $\mu$mol/L = 0.0113 mg/dL
\end{itemize}

\section*{Follow-up Tests}
\begin{itemize}
  \item \textbf{eGFR (automatically calculated):} Essential CKD staging. If eGFR <60, repeat in 3 months to confirm chronicity. If eGFR declines >5 mL/min/1.73 m\textsuperscript{2} per year, this represents accelerated progression.
  \item \textbf{BUN and BUN:Creatinine ratio:} Classify AKI etiology. Prerenal >20:1; intrinsic 10--15:1.
  \item \textbf{Cystatin C:} Useful alternative GFR marker in patients with abnormal muscle mass, cirrhosis, or when creatinine-based eGFR is unreliable. Cystatin C is produced by all nucleated cells and is independent of muscle mass. The combined creatinine-cystatin C CKD-EPI equation is more accurate than either alone.
  \item \textbf{Urinalysis with microscopy:} Muddy brown casts and RTECs $\rightarrow$ ATN. Dysmorphic RBCs and RBC casts $\rightarrow$ glomerulonephritis. WBC casts $\rightarrow$ interstitial nephritis.
  \item \textbf{Urine albumin-to-creatinine ratio (UACR) or urine protein-to-creatinine ratio (UPCR):} Essential for CKD staging and prognosis. Combined eGFR + albuminuria staging provides CKD risk stratification (KDIGO heat map).
  \item \textbf{Fractional excretion of sodium (FENa) and urea (FEUrea):} FENa <1\%, FEUrea <35\% $\rightarrow$ prerenal. FENa >2\%, FEUrea >50\% $\rightarrow$ ATN. FEUrea more reliable in patients on diuretics.
  \item \textbf{Renal ultrasound:} Kidney size, cortical thickness, echogenicity, cysts, hydronephrosis. Small kidneys (<8--9 cm) suggest chronicity.
  \item \textbf{Serologic workup for glomerulonephritis:} ANA, ANCA, anti-GBM, complement C3/C4, anti-dsDNA, anti-streptolysin O, cryoglobulins, hepatitis B/C, HIV.
  \item \textbf{Serum protein electrophoresis (SPEP) and serum free light chains:} For suspected multiple myeloma.
  \item \textbf{Kidney biopsy:} Definitive diagnosis for glomerular disease, unexplained AKI, or rapidly progressive GN.
\end{itemize}

\section*{References}
\bibliographystyle{unsrtnat}
\bibliography{references}

\clearpage
\section*{Regulatory Status and Authorization Date}
This chapter describes a test category, not a specific commercial product. FDA, CDSCO, EU IVDR, and MHRA approval or registration dates therefore cannot be assigned without the assay manufacturer, product name, instrument, and intended use. For laboratory-developed tests, an individual agency approval date may not exist. See the Preface for the interpretation of regulatory status.

\clearpage
